\documentclass{article}
\usepackage{anyfontsize}
\usepackage[UTF8]{ctex} % 如果需要支持中文
\usepackage{fontspec} 
\setCJKmainfont{SourceHanSansCN-Light}[
    BoldFont=SourceHanSansCN-Medium
]

\usepackage[a4paper, left=0.1in, right=0.1in, top=0.05in, bottom=0.05in]{geometry} % 设置四边页边距为0.1英寸{geometry} % 设置页边距为0.5英寸
\usepackage{enumitem} % 用于自定义列表
\usepackage{amsmath}  % 数学公式支持
\usepackage{tikz}
\usetikzlibrary{arrows.meta, positioning}
\setlength{\parindent}{0pt} % 不缩进
\usepackage{etoolbox} % 用于列表操作
%调整类代码
\usepackage{comment}
\usepackage{tocloft} % 用于定制目录 样式
\usepackage{hyperref} %不需要目录盒子注释这
\usepackage{ifthen}
\usepackage{tikz}
\usetikzlibrary{arrows.meta}
\usepackage{titletoc}
\usepackage{graphicx}
\usepackage{float}

\usepackage{amssymb}  % 提供数学符号，包括\therefore
%目录设定
%快捷键 CMD + / 注释
%  \renewcommand{\cftdot}{} % 隐藏目录中的页码
%  \cftpagenumbersoff{section}
%  \cftpagenumbersoff{subsection}
%  \makeatletter
%  \renewcommand{\@pnumwidth}{0em}  % 设置页码宽度为0
%  \renewcommand{\@tocrmarg}{0em}   % 设置右边距为0
%  \makeatother
 

\usepackage{environ}

\newif\ifshowcontent  % 定义一个条件判断变量
\showcontenttrue    % 设置为 false，隐藏所有 question 环境的内容

\newcounter{question} % 题目计数器
\setcounter{question}{0}
\NewEnviron{question}{%
    \ifshowcontent
        \par\stepcounter{question}\textbf{\thequestion.} \BODY\par % 如果显示内容，则输出
    \fi
}

\NewEnviron{choices}{%
    \ifshowcontent
        \begin{enumerate}[label=\Alph*., leftmargin=*]
            \BODY
        \end{enumerate}\par
    \fi
}

\NewEnviron{correctanswer}{%
    \ifshowcontent
        \textbf{答案:}\BODY\par
    \fi
}



\newenvironment{explanation}
{\textbf{解析:}\par}
{\par}



% 新增考察方面命令
\newcommand{\checklist}[1]{
    \textbf{考察:} #1 \par % 在题目下方显示考察方面
    \addcontentsline{toc}{subsection}{题号\thequestion 考察: #1} % 将考察内容添加到目录
    \vspace{2em} % 添加1em的垂直空间
}

\graphicspath{{images/}} %统一


\begin{document}
 %\fontsize{22}{31}\selectfont
 \tableofcontents % 添加目录
\newpage
%\pagestyle{empty}




\section{考点1:对冲基金}
\setcounter{question}{0}

\begin{question}
    When distinguishing between strategy risks and structural risks in a fund of hedge funds, which of the following would be considered strategy risk(s)?  
    \begin{itemize}
        \item I. Trading liquidity risk.  
        \item II. The extent and form of management oversight.  
        \item III. The risk of poor information reporting systems.  
        \item IV. The risk of high ownership concentration of hedge fund shares.
    \end{itemize}
\end{question}

\begin{choices}
    \item I, III, and IV only
    \item I only
    \item II and III only
    \item II and IV only
\end{choices}

\begin{correctanswer}
    B
\end{correctanswer}

\begin{explanation}
    \textbf{陈述I分析：策略风险}
    \begin{itemize}
        \item 交易流动性风险属于策略风险；它源于投资策略的执行过程，如持有非流动性资产、使用杠杆、快速交易等策略特征，是策略本身带来的风险。
    \end{itemize}

    \textbf{陈述II分析：结构性风险}
    \begin{itemize}
        \item 管理监督的范围和形式属于结构性风险；这涉及基金的治理结构、管理层的监督机制，属于基金的组织和运营结构问题，而非投资策略本身。
    \end{itemize}

    \textbf{陈述III分析：结构性风险}
    \begin{itemize}
        \item 信息报告系统不良的风险属于结构性风险；这涉及基金的运营基础设施、信息系统质量，属于运营风险，而非投资策略风险。
    \end{itemize}

    \textbf{陈述IV分析：结构性风险}
    \begin{itemize}
        \item 对冲基金份额高度集中的风险属于结构性风险；这涉及基金的投资者结构、所有权集中度，属于基金结构问题，可能影响基金治理和运营，而非策略风险。
    \end{itemize}
\end{explanation}
\checklist{对冲基金策略风险与结构性风险}


\begin{question}
    Hedge fund managers following a convertible arbitrage strategy are said to be:
\end{question}

\begin{choices}
    \item long gamma and short vega
    \item short gamma and short vega
    \item long gamma and long vega
    \item short gamma and long vega
\end{choices}

\begin{correctanswer}
    C
\end{correctanswer}

\begin{explanation}
    \textbf{理解可转债套利策略的风险敞口}
    \begin{itemize}
        \item 可转债套利策略通常买入可转债（包含嵌入看涨期权）并做空标的股票进行delta对冲。
    \end{itemize}

    \textbf{Gamma敞口分析：Long Gamma}
    \begin{itemize}
        \item 可转债中的嵌入期权是long gamma（持有期权）；虽然做空股票进行delta对冲，但由于delta会随股价变化而变化（delta对冲不完美），策略整体暴露于gamma风险，是long gamma。
    \end{itemize}

    \textbf{Vega敞口分析：Long Vega}
    \begin{itemize}
        \item 可转债中的嵌入期权是long vega（波动率上升对期权价值有利）；做空股票没有vega敞口；因此策略整体是long vega，暴露于标的股票波动率变化的风险。
    \end{itemize}
\end{explanation}
\checklist{可转债套利的希腊字母敞口}


\begin{question}
    The Quantum Global Fund is a global macro hedge fund designed to provide low correlations with European assets. Sarah Chen, a fund of hedge funds manager, is analyzing the Quantum Global Fund for signs of style drift. Chen notes the following findings: 
    \begin{itemize}
        \item I. The R² of the fund versus the global macro peer group has changed from 0.65 to 0.72 over the past 12 months.  
        \item II. Due to outstanding returns, assets in the fund have increased from €50 million to €350 million over the past 12 months.  
        \item III. The fund made a major shift in allocation by moving 35 percent of its holdings from Latin American equities to European equities.  
        \item IV. After a recent trip to Germany, the fund manager gained confidence in his existing German equity holdings and levered his existing 4\% weighting in Germany by a 8 to 1 ratio.  
    \end{itemize}
    Which of Chen's findings are indicators that the Quantum Global Fund is at risk for style drift?
\end{question}

\begin{choices}
    \item II and IV only
    \item I and II only
    \item II and III only
    \item I, III and IV only
\end{choices}

\begin{correctanswer}
    A
\end{correctanswer}

\begin{explanation}
    \textbf{I：不是风格漂移信号}
    \begin{itemize}
        \item R²从0.65→0.72说明与同类策略更一致，非偏离迹象。
    \end{itemize}

    \textbf{II：风格漂移风险信号}
    \begin{itemize}
        \item 规模从€50m→€350m，资产膨胀可能迫使策略、标的与交易方式改变（容量约束）。
    \end{itemize}

    \textbf{III：通常不构成风格漂移}
    \begin{itemize}
        \item 全球宏观框架内从拉美转向欧洲属于地域再配置，未必改变策略本质。
    \end{itemize}

    \textbf{IV：风格漂移风险信号}
    \begin{itemize}
        \item 将德国敞口加8:1杠杆，显著改变风险因子暴露与波动/尾部特征。
    \end{itemize}
\end{explanation}
\checklist{对冲基金风格漂移识别}


\begin{question}
    There has been a recent push for financial market participants to consider increasing the number of available reference rates, providing alternatives to the London Interbank Offered Rate (LIBOR) and Euribor. Which of the following criticisms of the commonly used reference rates has encouraged the need for alternative reference rates?
\end{question}

\begin{choices}
    \item Commonly used rates contain credit risk, making them unsuitable for some transactions.
    \item Overnight markets are less liquid, making them unappealing for use in swap contracts.
    \item Banks have become increasingly reliant on unsecured funding to better manage credit risks.
    \item LIBOR and Euribor are not consistently published on dedicated sites.
\end{choices}

\begin{correctanswer}
    A
\end{correctanswer}

\begin{explanation}
    \textbf{A选项正确。}  
    \begin{itemize}
        \item 常用参考利率包含信用风险，这使得它们不适合某些需要无风险利率的交易。金融危机后，银行风险分散化，使得包含信用风险的参考利率对某些交易不再适用。
    \end{itemize}

    \textbf{B选项错误。}  
    \begin{itemize}
        \item 隔夜市场实际上流动性更高，而不是更低。
    \end{itemize}

    \textbf{C选项错误。}  
    \begin{itemize}
        \item 银行越来越依赖批发融资来管理信用风险，而不是无担保融资。
    \end{itemize}

    \textbf{D选项错误。}  
    \begin{itemize}
        \item LIBOR和Euribor在专用网站上持续发布，这不是问题所在。
    \end{itemize}
\end{explanation}
\checklist{参考利率的信用风险问题}


\begin{question}
    A hedge fund database is constructed as follows. The first year is 2000. All funds existing at the end of 2000 that were willing to report their verified returns for that year are included. The database was extended by asking funds for verified returns before 2000. Subsequently, funds are added as they are willing to report verified returns. If a fund stops reporting, its returns are deleted from the database, but funds agree to keep reporting even if they stop accepting new investors. Consider the following statements:  
    \begin{itemize}
        \item I. The database suffers from backfilling bias.  
        \item II. The database suffers from survivorship bias.  
        \item III. The database suffers from an errors-in-variables bias.  
        \item IV. The equally-weighted annual return average will underestimate expected hedge fund performance.  
    \end{itemize}
    Which statements are correct?
\end{question}

\begin{choices}
    \item All the above statements are correct.
    \item Statements I and II are correct.
    \item Statements I, II, and III are correct.
    \item Statements II and IV are correct.
\end{choices}

\begin{correctanswer}
    B
\end{correctanswer}

\begin{explanation}
    \textbf{陈述I分析：正确}
    \begin{itemize}
        \item 数据库存在回填偏差（backfilling bias）；数据库要求基金提供2000年之前的历史数据，只有表现好的基金才愿意提供历史数据，导致历史数据偏向表现好的基金，产生回填偏差。
    \end{itemize}

    \textbf{陈述II分析：正确}
    \begin{itemize}
        \item 数据库存在生存偏差（survivorship bias）；如果基金停止报告，其收益会被删除；表现差的基金（可能关闭或停止报告）会被移除，只有持续报告的表现较好的基金保留在数据库中，导致生存偏差。
    \end{itemize}

    \textbf{陈述III分析：错误}
    \begin{itemize}
        \item 变量误差偏差（errors-in-variables bias）通常出现在回归分析中，当解释变量存在测量误差时；这不是数据库构建本身的问题，而是回归分析中的统计问题；数据库构建方式不直接导致变量误差偏差。
    \end{itemize}

    \textbf{陈述IV分析：错误}
    \begin{itemize}
        \item 等权重年度收益平均值会高估而非低估预期对冲基金表现；由于回填偏差（历史数据偏向表现好的基金）和生存偏差（表现差的基金被移除），平均收益会被高估，而非低估。
    \end{itemize}
\end{explanation}
\checklist{对冲基金数据库偏差类型}


\begin{question}
    The Alpha Hedge Fund has the following description of its activities. It uses simultaneous long and short positions in equities with a net beta close to zero. Which of the following statements about Alpha Hedge Fund are correct?  
    \begin{itemize}
    \item I. It uses a directional strategy.  
    \item II. It is a relative value hedge fund.  
    \item III. This fund is exposed to idiosyncratic risks.
    \end{itemize}
\end{question}

\begin{choices}
    \item I and II
    \item II and III
    \item I and III
    \item II only
\end{choices}

\begin{correctanswer}
    B
\end{correctanswer}

\begin{explanation}
    \textbf{陈述I分析：错误}
    \begin{itemize}
        \item 该基金不使用方向性策略；方向性策略通常有显著的市场敞口（beta不为零），通过预测市场方向获利；该基金净beta接近零，说明它消除了市场风险，不是方向性策略。
    \end{itemize}

    \textbf{陈述II分析：正确}
    \begin{itemize}
        \item 该基金是相对价值对冲基金；相对价值策略通过同时做多和做空相关证券来获取相对价值差异，通常beta接近零；该基金的特征（同时做多做空、净beta接近零）符合相对价值策略。
    \end{itemize}

    \textbf{陈述III分析：正确}
    \begin{itemize}
        \item 该基金暴露于特质风险（idiosyncratic risk）；虽然市场风险被对冲（beta接近零），但做多和做空的股票可能不完全相关，个股的特质风险仍然存在；如果做多股票表现差于做空股票，基金仍会遭受损失。
    \end{itemize}
\end{explanation}
\checklist{对冲基金策略类型识别}


\begin{question}
    Jennifer Lee manages a hedge fund for a mid-sized investment firm. The fund frequently changes styles according to identified profit opportunities. At the beginning of the year, the fund took a long position in 12-year subordinated 7\% coupon debt issued by a firm expected to undergo reorganization under Chapter 11. Lee felt that analysts had been paying too little attention to the issuer. Six months later, the fund completed a second transaction involving a long position in Euro and a short position in British Pounds based on forecasted movements in interest rates in the two regions. What two hedge fund strategies are most likely being employed by Lee's hedge fund?
\end{question}

\begin{choices}
    \item Distressed securities strategy and equity long/short strategy.
    \item Fixed-income arbitrage and global macro strategy.
    \item Distressed securities strategy and global macro strategy.
    \item Fixed-income arbitrage and equity long/short strategy.
\end{choices}

\begin{correctanswer}
    C
\end{correctanswer}

\begin{explanation}
    \textbf{交易1判定：困境证券（Distressed Securities）}
    \begin{itemize}
        \item 做多12年期次级7\%票息债，发行人预计进入Chapter 11重组，典型困境债投资。
    \end{itemize}

    \textbf{交易2判定：全球宏观（Global Macro）}
    \begin{itemize}
        \item 做多欧元、做空英镑，基于两地利率预期差的宏观自上而下货币交易，属全球宏观。
    \end{itemize}
    \textbf{选项逐一：}
    \begin{itemize}
        \item A 错：第二笔不是股票多空。
        \item B 错：第一笔不是固定收益套利，而是困境债。
        \item C 对：困境证券 + 全球宏观。
        \item D 错：两笔均非固收套利或股多空。
    \end{itemize}
\end{explanation}
\checklist{对冲基金策略类型识别}


\begin{question}
    Michael Chen, FRM, is a credit analyst. In a presentation to investors on the recent Asian financial crisis, Chen mentioned that during the first stages of the crisis, there was generally a positive correlation between bank and sovereign credit default swap (CDS) spreads as both sovereign and bank CDS spreads rose dramatically. However, as the Asian banking sector problems became much more evident, the correlation turned negative as bank CDS spreads rose while sovereign spreads declined. Are Chen's statements with regard to positive and negative correlation correct?
\end{question}

\begin{choices}
    \item Yes, both statements are correct.
    \item Yes, but only the first statement is correct.
    \item No, but only the second statement is correct.
    \item No, neither statement is correct.
\end{choices}

\begin{correctanswer}
    D
\end{correctanswer}

\begin{explanation}
    \textbf{陈述I：错误}
    \begin{itemize}
        \item 危机初期银行与主权CDS利差的相关性并非“通常为正”；受政府隐性担保、救助预期与市场传染等影响，相关性不稳定，难以一概而论。
    \end{itemize}

    \textbf{陈述II：错误}
    \begin{itemize}
        \item 银行业问题暴露后，常见的是“厄运循环”（doom loop）：银行与主权风险相互强化，利差多呈同向上升，相关性趋于正而非负。
    \end{itemize}

\end{explanation}
\checklist{CDS利差相关性变化}


\begin{question}
    Fixed income arbitrage funds attempt to obtain profits by exploiting inefficiencies and price anomalies between related fixed-income securities. The fund managers try to limit volatility by hedging exposure to interest rate risk. Which of the following types of fixed-income trades bets that the fixed side of a spread will stay higher than the floating side of a spread?
\end{question}

\begin{choices}
    \item Swap spread trade.
    \item Credit arbitrage trades.
    \item Mortgage spread trades.
    \item Fixed-income volatility trades.
\end{choices}

\begin{correctanswer}
    A
\end{correctanswer}

\begin{explanation}
    \textbf{A选项正确。}  
    \begin{itemize}
    \item 互换利差交易是对固定端利差将保持在浮动端利差之上的押注，并预期利差保持在历史趋势的合理范围内。
    \end{itemize}
    \textbf{B选项错误。}  
    \begin{itemize}
        \item 信用套利交易主要关注信用利差，不是固定端与浮动端的利差。
    \end{itemize}
    \textbf{C选项错误。}  
    \begin{itemize}
    \item 抵押贷款利差交易主要关注抵押贷款证券的利差，不是固定端与浮动端的利差。
    \end{itemize}
    \textbf{D选项错误。} 
    \begin{itemize} 
    \item 固定收益波动率交易关注的是波动率，不是利差的固定端与浮动端。
    \end{itemize}
\end{explanation}
\checklist{固定收益套利策略类型}


\begin{question}
    Quantum Tech (QT) and Nexus Systems (NS) are both leading semiconductor manufacturers. Earlier this year, both companies introduced new AI chips that are comparable in almost every category. However, after both companies release pricing for their new chips, NS's model is 25\% less expensive than QT's. As a result, QT's stock price declined while NS's stock price rose. Subsequently, QT and NS announce they have entered into merger discussions where the terms would give QT shareholders 1 share of NS per 2 shares of QT previously held. Post the announcement, QT's stock is trading at USD 30 and NS's stock is trading at USD 65. If you are confident that the merger will be completed, assuming zero transaction costs, which investment should you make?
\end{question}

\begin{choices}
    \item Buy 200 shares of QT and short 100 shares of NS.
    \item Short 200 shares of QT and buy 100 shares of NS.
    \item Buy 200 shares of QT and buy 100 shares of NS.
    \item Short 200 shares of QT and short 100 shares of NS.
\end{choices}

\begin{correctanswer}
    B
\end{correctanswer}

\begin{explanation}
  
    \textbf{思路：合并换股比与当前价格比对比}
    \begin{itemize}
        \item 换股比：2 QT = 1 NS（合并后每2股QT换1股NS，套保比2:1）。
        \item 现价比：NS/QT = 65/30 = 2.17 > 2（NS相对便宜，QT相对贵）。
    \end{itemize}

    \textbf{操作与利润}
    \begin{itemize}
        \item 做空200股QT（+\$6{,}000），买入100股NS（-\$6{,}500），净流出\$500。
        \item 合并完成：空头200 QT以换股获得的100 NS对冲（头寸抵消），锁定价差损益=\$65×100 - \$30×200 = -\$500；与期初净流出相抵，利润=\$0? 
        \item 等价看法：价差（理论2 vs 实际2.17）每套2:1头寸的错定价为\$65 - 2×\$30 = \$5；做1套需净投入\$500（买1 NS、卖空2 QT），合并时收敛，收益\$5×100=\$500，对应B的组合。
    \end{itemize}

   
\end{explanation}
\checklist{合并套利策略}



\begin{question}
    Sarah Kim, FRM, is considering investing a client into distressed hedge funds. Which of the following investments would serve as the best proxy for the types of returns to expect?
\end{question}

\begin{choices}
    \item Convertible bonds.
    \item Small-cap equities.
    \item Managed futures.
    \item High-yield bonds.
\end{choices}

\begin{correctanswer}
    D
\end{correctanswer}

\begin{explanation}
    \textbf{A选项错误。}  
    \begin{itemize}
        \item 可转债虽然包含信用风险，但主要风险来自股票价格变动，不是困境证券策略的最佳代理。
    \end{itemize}

    \textbf{B选项错误。}  
    \begin{itemize}
        \item 小盘股主要暴露于市场风险和流动性风险，不是信用风险。
    \end{itemize}

    \textbf{C选项错误。}  
    \begin{itemize}
        \item 管理期货主要暴露于商品和金融期货价格变动，与困境证券策略无关。
    \end{itemize}

    \textbf{D选项正确。}  
    \begin{itemize}
        \item 困境证券对冲基金主要暴露于低信用评级公司的信用风险，公开交易的高收益债券是预期回报的最佳代理。
    \end{itemize}
\end{explanation}
\checklist{困境证券策略的回报代理}



\begin{question}
    Identify the risks in a fixed-income arbitrage strategy that takes long positions in interest rate swaps hedged with short positions in Treasuries.
\end{question}

\begin{choices}
    \item The strategy could lose from decreases in the swap-Treasury spread.
    \item The strategy could lose from increases in the Treasury rate, all else fixed.
    \item The payoff in the strategy has negative skewness.
    \item The payoff in the strategy has positive skewness.
\end{choices}

\begin{correctanswer}
    C
\end{correctanswer}

\begin{explanation}
    \textbf{A选项错误。}  
    \begin{itemize}
        \item 该策略应该从互换-国债利差的下降中获益，而不是损失。
    \end{itemize}

    \textbf{B选项错误。}  
    \begin{itemize}
        \item 该策略应该从国债利率的上升中获益，而不是损失。
    \end{itemize}

    \textbf{C选项正确。}  
    \begin{itemize}
        \item 该策略的收益分布取决于互换-国债利差的分布。由于利差不能低于零，上行空间有限，因此收益分布呈现负偏态。
    \end{itemize}

    \textbf{D选项错误。}  
    \begin{itemize}
        \item 该策略的收益分布不是正偏态。
    \end{itemize}
\end{explanation}
\checklist{固定收益套利策略的风险特征}


\begin{question}
    What critical shift occurred in the hedge fund industry following the collapse of Long-Term Capital Management (LTCM) in 1998 and the dot-com bubble burst in 2001?
\end{question}

\begin{choices}
    \item There was a significant drop in assets under management in the hedge fund industry.
    \item There was a large influx of institutional investors investing in hedge funds.
    \item Reporting within the hedge fund industry became more regulated than mutual funds.
    \item There was a significant increase in hedge fund failures.
\end{choices}

\begin{correctanswer}
    B
\end{correctanswer}

\begin{explanation}
    \textbf{A选项错误。}  
    \begin{itemize}
        \item 对冲基金行业的管理资产实际上增加了，而不是下降。
    \end{itemize}

    \textbf{B选项正确。}  
    \begin{itemize}
        \item 在互联网泡沫破裂后，对冲基金以较低的标准差跑赢了标普500指数，这吸引了机构投资者的涌入。
    \end{itemize}

    \textbf{C选项错误。}  
    \begin{itemize}
        \item 对冲基金行业的报告要求并没有变得比共同基金更严格。
    \end{itemize}

    \textbf{D选项错误。}  
    \begin{itemize}
        \item 对冲基金失败率并没有显著增加。
    \end{itemize}
\end{explanation}
\checklist{对冲基金行业的历史转折点}


\begin{question}
    What would be an ideal approach for a hedge fund investor who is concerned about the problem of risk sharing asymmetry between principals and agents within the hedge fund industry?
\end{question}

\begin{choices}
    \item Focus on investing in funds for which the fund managers have a good portion of their own wealth invested.
    \item Focus on diversifying among the various niche hedge fund strategies.
    \item Focus on funds with improved operational efficiency and transparent corporate governance.
    \item Focus on large funds from the "foresight-assisted" group.
\end{choices}

\begin{correctanswer}
    A
\end{correctanswer}

\begin{explanation}
    \textbf{A选项正确。}  
    \begin{itemize}
        \item 对冲基金行业的激励费结构多年来没有真正改变，基金经理有动机为了赚取费用而承担过度风险。因此，投资于基金经理将大量自有财富投入的基金是最理想的方法，这样可以确保基金经理与投资者利益一致。
    \end{itemize}

    \textbf{B选项错误。}  
    \begin{itemize}
        \item 分散投资于不同策略并不能解决委托代理问题。
    \end{itemize}

    \textbf{C选项错误。}  
    \begin{itemize}
        \item 虽然运营效率和公司治理很重要，但不能直接解决风险分担不对称问题。
    \end{itemize}

    \textbf{D选项错误。}  
    \begin{itemize}
        \item 基金规模大小与解决委托代理问题无关。
    \end{itemize}
\end{explanation}
\checklist{对冲基金委托代理问题的解决方案}



\begin{question}
    Identify the risks in a convertible arbitrage strategy that takes long positions in convertible bonds hedged with short positions in Treasuries and the underlying stock.
\end{question}

\begin{choices}
    \item Short implied volatility
    \item Long duration
    \item Long stock delta
    \item Positive gamma
\end{choices}

\begin{correctanswer}
    D
\end{correctanswer}

\begin{explanation}
    \textbf{A选项错误。}  
    \begin{itemize}
        \item 该策略做多可转债（包含转换期权），因此是long implied volatility，不是short。
    \end{itemize}

    \textbf{B选项错误。}  
    \begin{itemize}
        \item 该策略通过做空国债对冲了利率风险，因此不是long duration。
    \end{itemize}

    \textbf{C选项错误。}  
    \begin{itemize}
        \item 该策略通过做空标的股票对冲了方向性风险，因此delta接近零，不是long delta。
    \end{itemize}

    \textbf{D选项正确。}  
    \begin{itemize}
        \item 做多期权（可转债的转换期权）具有正gamma，这是该策略的主要风险敞口。
    \end{itemize}
\end{explanation}
\checklist{可转债套利策略的风险敞口}




\begin{question}
    Which of the following hedge fund strategies would be characterized as an "asset allocation" strategy that performs best during extreme moves in the currency markets?
\end{question}

\begin{choices}
    \item Global macro.
    \item Risk arbitrage.
    \item Dedicated short bias.
    \item Long/short equity.
\end{choices}

\begin{correctanswer}
    A
\end{correctanswer}

\begin{explanation}
    \textbf{A选项正确。}  
    \begin{itemize}
        \item 全球宏观基金在货币市场出现极端波动时表现最佳。全球宏观是一种资产配置策略，基金经理在不同市场中进行机会性投资，与股票市场相关性较低。
    \end{itemize}

    \textbf{B选项错误。}  
    \begin{itemize}
        \item 风险套利主要关注并购交易，不是资产配置策略。
    \end{itemize}

    \textbf{C选项错误。}  
    \begin{itemize}
        \item 专门做空策略主要关注股票市场，不是资产配置策略。
    \end{itemize}

    \textbf{D选项错误。}  
    \begin{itemize}
        \item 股票多空策略主要关注股票市场，不是资产配置策略。
    \end{itemize}
\end{explanation}
\checklist{全球宏观策略的特征}


\begin{question}
    When comparing hedge fund performance between the 2002-2010 period and earlier time periods, what was the trend in monthly alpha for large hedge funds?
\end{question}

\begin{choices}
    \item Alpha showed an increase during the 2002-2010 period.
    \item Alpha remained stable across both time periods.
    \item A "foresight-assisted" portfolio exhibited no statistically significant alpha in the 2002-2010 period.
    \item Alpha experienced a decline during the 2002-2010 period.
\end{choices}

\begin{correctanswer}
    D
\end{correctanswer}

\begin{explanation}
    \textbf{A选项错误。}
    \begin{itemize}
        \item 2002-2010年期间，对冲基金行业竞争加剧，alpha并未上升
    \end{itemize}
    
    \textbf{B选项错误。}
    \begin{itemize}
        \item alpha在不同时期并非保持稳定，而是呈现下降趋势
    \end{itemize}
    
    \textbf{C选项错误。}
    \begin{itemize}
        \item "前瞻性辅助"投资组合的alpha显著性检验结果与整体alpha下降趋势无关
    \end{itemize}
    
    \textbf{D选项正确。}
    \begin{itemize}
        \item 由于对冲基金行业竞争加剧，2002-2010年期间大型对冲基金的月度alpha出现下降
    \end{itemize}
\end{explanation}

\checklist{对冲基金alpha下降趋势}


\begin{question}
    An investment manager is presenting hedge fund strategies to junior analysts. Which statement about hedge fund strategies and their risk-return characteristics would be accurate?
\end{question}

\begin{choices}
    \item Merger arbitrage funds typically carry no idiosyncratic risk, and any such risk cannot be diversified away.
    \item Managed futures funds using trend-following strategies focus on single asset classes and employ lower leverage than other strategies.
    \item Distressed hedge funds hold highly illiquid securities, with tail risk primarily from large movements in short-term interest rates.
    \item Global macro funds use relatively passive strategies where managers adjust positions only during major macroeconomic developments.
\end{choices}

\begin{correctanswer}
    C
\end{correctanswer}

\begin{explanation}
    \textbf{A选项错误。}
    \begin{itemize}
        \item 合并套利基金确实存在特质风险，且这种风险可以通过分散化来降低
    \end{itemize}
    
    \textbf{B选项错误。}
    \begin{itemize}
        \item 管理期货基金通常跨多个资产类别，且使用较高杠杆
    \end{itemize}
    
    \textbf{C选项正确。}
    \begin{itemize}
        \item 困境证券基金持有高度非流动性证券，尾部风险主要来自短期利率的大幅波动
    \end{itemize}
    
    \textbf{D选项错误。}
    \begin{itemize}
        \item 全球宏观基金采用主动策略，而非被动策略
    \end{itemize}
\end{explanation}

\checklist{对冲基金策略特征}


\begin{question}
    A portfolio manager at a hedge fund is explaining portfolio construction to new analysts. When discussing alpha refinement, which statement correctly describes the primary goal and mechanics of refining alphas?
\end{question}

\begin{choices}
    \item Minimize noise in portfolio construction techniques like quadratic programming by adjusting alphas after construction.
    \item Simplify portfolio construction by preprocessing alphas to reflect the manager's constraints and objectives.
    \item Increase robustness of portfolio construction by making subtle alpha adjustments to facilitate computer model inputs.
    \item Rebalance portfolios using alphas updated with new information about portfolio and benchmark securities.
\end{choices}

\begin{correctanswer}
    B
\end{correctanswer}

\begin{explanation}
    \textbf{A选项错误。}
    \begin{itemize}
        \item alpha精炼不是在构建后调整，而是在构建前预处理
    \end{itemize}
    
    \textbf{B选项正确。}
    \begin{itemize}
        \item alpha精炼的主要目标是通过预处理alpha来简化投资组合构建，使其反映管理者的约束和目标
    \end{itemize}
    
    \textbf{C选项错误。}
    \begin{itemize}
        \item 精炼alpha不是为了增加技术稳健性，而是为了简化构建过程
    \end{itemize}
    
    \textbf{D选项错误。}
    \begin{itemize}
        \item 这不是alpha精炼的主要目标，而是投资组合再平衡的目的
    \end{itemize}
\end{explanation}

\checklist{alpha精炼目标}




\newpage



\section{考点2:对特定管理人和基金进行尽职调查}
\setcounter{question}{0}

\begin{question}
    When conducting due diligence on a potential investment manager, which of the following factors should be considered least important by the investor?
\end{question}

\begin{choices}
    \item Risk controls.
    \item Business model.
    \item Past performance.
    \item Investment process.
\end{choices}

\begin{correctanswer}
    C
\end{correctanswer}

\begin{explanation}
    \textbf{A选项错误。}
    \begin{itemize}
        \item 风险控制是尽职调查的核心要素，对评估投资经理至关重要
    \end{itemize}
    
    \textbf{B选项错误。}
    \begin{itemize}
        \item 商业模式是评估投资经理可持续性和稳定性的重要因素
    \end{itemize}
    
    \textbf{C选项正确。}
    \begin{itemize}
        \item 过去表现虽然重要，但不应过度关注，因为历史业绩不能保证未来表现。风险控制、商业模式和投资流程才是尽职调查的基本组成部分
    \end{itemize}
    
    \textbf{D选项错误。}
    \begin{itemize}
        \item 投资流程是评估投资经理方法论和一致性的关键要素
    \end{itemize}
\end{explanation}

\checklist{尽职调查要素重要性}


\begin{question}
    A due diligence specialist is evaluating a hedge fund's risk management process for potential investment. Which statement best describes the evaluation criteria?
\end{question}

\begin{choices}
    \item Due to the critical importance of tail risk, the company should only invest if the fund fully accounts for fat tails using extreme value theory at the 99.99\% VaR level.
    \item Current best practices require funds to employ independent risk service providers who play important roles in risk-related decisions.
    \item When evaluating a leveraged fund, the specialist should assess how the fund estimates leverage-related risks, including funding liquidity risks during market stress periods.
    \item It is essential to evaluate the fund's valuation policy, and if more than 10\% of asset prices are based on model prices or broker quotes, investment should be avoided regardless of other fund information.
\end{choices}

\begin{correctanswer}
    C
\end{correctanswer}

\begin{explanation}
    \textbf{A选项错误。}
    \begin{itemize}
        \item 虽然尾部风险重要，但投资者应自行决定是否接受未对冲的风险或自行对冲
    \end{itemize}
    
    \textbf{B选项错误。}
    \begin{itemize}
        \item 独立风险服务提供商主要向投资者报告风险，不参与风险相关决策
    \end{itemize}
    
    \textbf{C选项正确。}
    \begin{itemize}
        \item 评估杠杆基金时，应重点考虑杠杆相关风险，特别是市场压力时期的融资流动性风险
    \end{itemize}
    
    \textbf{D选项错误。}
    \begin{itemize}
        \item 可接受的模型定价比例取决于基金策略，不能一概而论
    \end{itemize}
\end{explanation}

\checklist{风险管理评估标准}



\begin{question}
    Lisa Tahara, FRM, is evaluating an institutional investment in a hedge fund with volatile but generally positive historical returns. Which consideration about the fund's track record is least relevant for her investment decision?
\end{question}

\begin{choices}
    \item Size of investment assets.
    \item Absolute level of past returns.
    \item Verification of returns by a third party.
    \item Employment continuity of the investment team.
\end{choices}

\begin{correctanswer}
    B
\end{correctanswer}

\begin{explanation}
    \textbf{A选项错误。}
    \begin{itemize}
        \item 投资资产规模是评估基金稳定性和流动性的重要因素
    \end{itemize}
    
    \textbf{B选项正确。}
    \begin{itemize}
        \item 考虑到历史回报的波动性，绝对回报水平最不相关。过去回报不能保证未来类似表现，相对回报水平比绝对水平更重要
    \end{itemize}
    
    \textbf{C选项错误。}
    \begin{itemize}
        \item 第三方验证回报提供了计算公平性和准确性的保证
    \end{itemize}
    
    \textbf{D选项错误。}
    \begin{itemize}
        \item 投资团队的就业连续性有助于评估未来产生类似回报的可能性
    \end{itemize}
\end{explanation}

\checklist{投资决策相关因素}


\begin{question}
    Based on historical evidence, which of the following factors is least likely to cause the eventual failure of a hedge fund?
\end{question}

\begin{choices}
    \item Excessive controls in place.
    \item Taking on more systematic risk.
    \item Making decisions in a committee setting.
    \item Materially misstated financial statements.
\end{choices}

\begin{correctanswer}
    B
\end{correctanswer}

\begin{explanation}
    \textbf{A选项错误。}
    \begin{itemize}
        \item 过度控制可能导致费用过高和回报不足，从而可能导致基金失败
    \end{itemize}
    
    \textbf{B选项正确。}
    \begin{itemize}
        \item 承担更多系统性风险（即常规市场风险）不太可能导致失败，除非出现重大市场下跌。承担更多非系统性风险才更容易导致失败
    \end{itemize}
    
    \textbf{C选项错误。}
    \begin{itemize}
        \item 委员会决策过程中可能存在主导成员或成员不敢表达有效关切，增加失败风险
    \end{itemize}
    
    \textbf{D选项错误。}
    \begin{itemize}
        \item 重大财务造假是会计欺诈，显著增加基金最终失败的风险
    \end{itemize}
\end{explanation}

\checklist{对冲基金失败因素}


\begin{question}
    An insurance company's investment team is conducting due diligence on a hedge fund. After comprehensive evaluation of past and current performance and management interviews, which finding would provide the strongest motivation to invest in the hedge fund?
\end{question}

\begin{choices}
    \item The hedge fund management possesses a strong combination of investment, operational, and business skills.
    \item The hedge fund's investment strategy has performed very well in the past.
    \item The hedge fund uses higher leverage and has a lower liquidity profile compared to most other hedge funds.
    \item The CEO of the hedge fund has substantial experience in investment strategy.
\end{choices}

\begin{correctanswer}
    A
\end{correctanswer}

\begin{explanation}
    \textbf{A选项正确。}
    \begin{itemize}
        \item 尽职调查过程分为四个关键领域：投资流程、风险管理、运营环境和商业模式评估。管理团队在投资、运营和商业技能方面的综合实力是最重要的考虑因素
    \end{itemize}
    
    \textbf{B选项错误。}
    \begin{itemize}
        \item 过去表现良好不能保证未来表现，不是最重要的投资动机
    \end{itemize}
    
    \textbf{C选项错误。}
    \begin{itemize}
        \item 高杠杆和低流动性特征会增加风险，不是积极的投资因素
    \end{itemize}
    
    \textbf{D选项错误。}
    \begin{itemize}
        \item 仅CEO有投资策略经验不足以证明整个管理团队的综合能力
    \end{itemize}
\end{explanation}

\checklist{尽职调查关键因素}



\begin{question}
    Senior analysts at an investment firm are meeting with managers of a small hedge fund for due diligence. After reviewing the investment process and risk controls, they focus on operational environment and business model. Which situation would constitute a red flag from an operational perspective?
\end{question}

\begin{choices}
    \item The hedge fund's financial statements are unqualified.
    \item The hedge fund partnership agreement contains an indemnification provision.
    \item The hedge fund does not use a compliance service provider.
    \item The hedge fund's general partner has made significant withdrawals.
\end{choices}

\begin{correctanswer}
    D
\end{correctanswer}

\begin{explanation}
    \textbf{A选项错误。}
    \begin{itemize}
        \item 无保留意见的财务报表是正常且积极的情况
    \end{itemize}
    
    \textbf{B选项错误。}
    \begin{itemize}
        \item 合伙协议中的赔偿条款是常见的法律保护措施
    \end{itemize}
    
    \textbf{C选项错误。}
    \begin{itemize}
        \item 不使用合规服务提供商可能是成本考虑，不一定是危险信号
    \end{itemize}
    
    \textbf{D选项正确。}
    \begin{itemize}
        \item 普通合伙人大量撤资可能表明对基金前景缺乏信心，是重要的运营风险信号
    \end{itemize}
\end{explanation}

\checklist{运营风险信号}



\newpage


\section{考点3:侦查投资经理的欺诈行为}
\setcounter{question}{0}

\begin{question}
    A junior examiner at a financial supervisory authority is studying US investment fraud cases, analyzing the connection between SEC financial disclosure filings (particularly Form ADV) and fraud instances. Based on the examiner's findings, which statement is correct?
\end{question}

\begin{choices}
    \item Disclosures of past conflicts of interest and prohibitions of future conflicts have proven effective in deterring fraud.
    \item Investors are typically compensated for fraud risk exposure, which decreases the quality of information disclosures by investment advisors.
    \item The benefits of using Form ADV filings in fraud protection models have always outweighed the marginal costs of collecting them.
    \item Using backward-looking fraud prediction models that relate past fraud with contemporaneous data exposes investors to survivorship bias.
\end{choices}

\begin{correctanswer}
    D
\end{correctanswer}

\begin{explanation}
    \textbf{A选项错误。}
    \begin{itemize}
        \item 过去利益冲突披露和未来冲突禁止并未被证明能有效阻止欺诈
    \end{itemize}
    
    \textbf{B选项错误。}
    \begin{itemize}
        \item 投资者通常不会因欺诈风险暴露而获得补偿，这也不会降低信息披露质量
    \end{itemize}
    
    \textbf{C选项错误。}
    \begin{itemize}
        \item Form ADV文件在欺诈保护模型中的收益并不总是超过收集的边际成本
    \end{itemize}
    
    \textbf{D选项正确。}
    \begin{itemize}
        \item 使用将过去欺诈与同期数据关联的回顾性欺诈预测模型会使投资者面临生存偏差
    \end{itemize}
\end{explanation}

\checklist{欺诈预测模型偏差}




\newpage


\end{document}